\begin{abstract}
마스킹 이산 확산 언어 모델은 범주형 잡음 제거 전이를 연속적으로 수행하여 텍스트를 생성한다.
본 연구는 토큰별 범주형 예측을 Fisher--Rao 통계 다양체 위의 점으로 취급하는 추론 시점 수정법을
다룬다. 제곱근 사상은 범주형 단체를 초구의 양의 직교영역에 매장하며, 이 공간에서는 닫힌형식의
구면 보간을 사용할 수 있다. 각 역방향 단계에서 이전 외삽 예측에서 현재 모델 예측을 통과하는
대원의 호를 재귀적으로 연장하고, 그 결과를 다시 범주형 분포로 변환한 뒤 원래 MDLM 또는 SEDD
전이에 삽입한다. 이 방법을 \textbf{측지 모멘텀 샘플링}이라 부른다.

정확한 기하학적 유도를 제시하고, 사영 이전 갱신이 지수사상에 의한 연장과 동치임을 증명하며,
작은 각도에서의 계수가 가변 스텝 2단계 Adams--Bashforth 외삽 계수와 일치함을 보인다. 동시에
이 유사성의 한계를 분명히 한다. 본 방법은 접벡터장이 아니라 예측 분포를 외삽하고, 확률적 범주형
점프를 유지하며, 양의 직교영역으로의 사영을 사용한다. 따라서 이 방법은 2차 정확도의 확률흐름
ODE 솔버가 아니며, 2차 수렴성을 주장하지 않는다. OpenWebText 체크포인트로 수행되어 보존된
실험 결과는 MDLM의 10--100 역방향 스텝과 SEDD의 10--50 스텝에서 GPT-2 생성 퍼플렉서티가
개선될 가능성을 보여준다.
\end{abstract}

\section{서론}

이산 확산 언어 모델은 병렬적이고 양방향적인 생성을 제공하지만, 정확한 샘플링에는 흔히 많은
역방향 전이가 필요하다. MDLM은 치환 매개화를 통해 흡수 상태 확산을 단순화하고 효율적인 캐시
샘플러를 제공한다 \citep{sahoo2024mdlm}. 반면 SEDD는 구체적인 스코어 비율을 학습하고 해석적인
역방향 전이를 허용한다 \citep{lou2024sedd}. 두 모델은 모두 어휘집 위의 범주형 분포를 반복적으로
생성하는 방식으로 작동한다.

정보기하는 이러한 분포에 자연스러운 기하를 부여한다. 범주형 단체의 내부에는 Fisher--Rao 계량이
정의되며 \citep{amari2016information}, 제곱근 사상은 상수 배율을 제외하면 이를 구면의 양의
직교영역과 동일시한다. 이 구조는 Fisher-Flow \citep{davis2024fisherflow}와 연속 리만 확산 언어
모델(RDLM) \citep{jo2025rdlm}의 토대이기도 하다. 이들 방법은 구면 위의 연속 흐름을 학습하거나
시뮬레이션한다. 본 연구의 설정은 다르다. 사전학습된 이산 MDLM과 SEDD는 그대로 유지하고,
추론 시점의 범주형 예측만 수정한다.

본 연구의 주요 기여는 다음과 같다.
\begin{itemize}
  \item 정확한 제곱근 매장과 구면 선형 보간(SLERP)을 사용해 범주형 예측을 재귀적으로 측지
  외삽하는 방법을 정식화한다.
  \item 사영 이전 갱신이 지수사상에 의한 연장임을 보이고, 국소적인 Adams--Bashforth 유사 계수
  구조를 유도한다.
  \item 동일한 기하 연산자를 MDLM의 정답 토큰 예측과 SEDD의 해석적 전이 분포라는 서로 다른
  두 인터페이스에 적용한다.
  \item 불변량과 고차 방법과의 유사성이 성립하는 정확한 범위를 명시한다. 특히 본
  알고리즘을 RDLM 확률흐름 ODE 솔버와 동일시하지 않는다.
  \item 여섯 가지 역방향 스텝 수에서 MDLM 및 SEDD의 품질 결과를 보고한다.
\end{itemize}
